%==============================================================================
%  ECON265 / CES365 -- Introduction to Econometrics
%  Reusable syllabus template.  Compile with:  pdflatex syllabus.tex  (x2)
%
%  >>> TO REUSE FOR A NEW SEMESTER, EDIT ONLY THE "COURSE CONFIGURATION"
%  >>> BLOCK BELOW.  The body text pulls everything from these macros.
%==============================================================================
\documentclass[11pt]{article}

%------------------------------------------------------------------------------
%  COURSE CONFIGURATION  -- edit me each semester
%------------------------------------------------------------------------------
\newcommand{\CourseNumber}{ECON265/CES365}
\newcommand{\CourseTitle}{Introduction to Econometrics}
\newcommand{\Semester}{Fall 2026}

% ---- Instructor ----
\newcommand{\InstructorName}{Moshi Alam}
\newcommand{\InstructorShort}{Moshi}
\newcommand{\InstructorEmail}{mdalam@clarku.edu}
\newcommand{\LectureRoom}{Sackler S311}
\newcommand{\LectureTimes}{Tue \& Thu, 10:25--11:40AM}
\newcommand{\OfficeRoom}{JC222}
\newcommand{\OfficeHours}{Thu 12:05--1:05PM, or by appointment}

% ---- Teaching assistant ----
\newcommand{\TAName}{Yujie Zheng}
\newcommand{\TAShort}{Yujie}
\newcommand{\TAthey}{she}                            % pronouns, subject
\newcommand{\TAthem}{her}                            % pronouns, object
\newcommand{\TAtheir}{her}                           % pronouns, possessive
\newcommand{\TAEmail}{YujiZeng@clarku.edu}
% Discussion sections: one \SectionLine{day+time}{room} per section
\newcommand{\TASessions}{%
  \SectionLine{Mon 10:00--10:50AM}{Carlson CH128}%
  \SectionLine{Wed 10:00--10:50AM}{Carlson CH120}%
}
\newcommand{\TAOfficeHours}{Mon \& Wed 11:00AM--12:00PM}
\newcommand{\TAOfficeRoom}{JC201}

% ---- Exams (set \ShowExamDates to 0 to hide the key-dates table) ----
\newcommand{\ShowExamDates}{1}
\newcommand{\MidtermOneDate}{Thursday, October 8, 2026}
\newcommand{\MidtermTwoDate}{Thursday, November 5, 2026}
\newcommand{\FinalDate}{Thursday, December 3, 2026}

% ---- Grade weights (must sum to 100) ----
\newcommand{\WtProblemSets}{10}
\newcommand{\WtQuizzes}{15}
\newcommand{\WtClassAtt}{10}
\newcommand{\WtDiscAtt}{5}
\newcommand{\WtMidtermOne}{20}
\newcommand{\WtMidtermTwo}{20}
\newcommand{\WtFinal}{20}

% ---- Assignment counts ----
\newcommand{\NumProblemSets}{5}
\newcommand{\NumQuizzes}{4}
\newcommand{\NumQuizzesCounted}{3}

%------------------------------------------------------------------------------
%  PREAMBLE  -- you should not need to touch anything below
%------------------------------------------------------------------------------
\usepackage[T1]{fontenc}
\usepackage[utf8]{inputenc}
\usepackage{newtxtext,newtxmath}
\usepackage{microtype}
\usepackage[margin=1in,top=0.9in,bottom=1in]{geometry}
\usepackage{enumitem}
\usepackage{booktabs}
\usepackage{tabularx}
\usepackage{xcolor}
\usepackage{titlesec}
\usepackage{fancyhdr}
\usepackage[hidelinks,breaklinks]{hyperref}

\definecolor{accent}{HTML}{7C2128}
\hypersetup{colorlinks=true,linkcolor=accent,urlcolor=accent,
  pdftitle={\CourseNumber: \CourseTitle{} (\Semester)},
  pdfauthor={\InstructorName}}

\titleformat{\section}{\normalfont\large\bfseries\color{accent}}{}{0pt}{}
\titlespacing*{\section}{0pt}{1.4ex plus .3ex}{0.6ex}
\titleformat{\subsection}[runin]{\normalfont\bfseries\itshape}{}{0pt}{}[:\hspace{0.4em}]
\titlespacing*{\subsection}{0pt}{1.0ex plus .2ex}{0.4em}

\setlist[itemize]{leftmargin=1.4em,topsep=2pt,itemsep=1.5pt,parsep=0pt}
\setlist[enumerate]{leftmargin=1.6em,topsep=2pt,itemsep=1.5pt,parsep=0pt}
\setlength{\parindent}{0pt}
\setlength{\parskip}{0.7ex plus .1ex}

\pagestyle{fancy}
\fancyhf{}
\renewcommand{\headrulewidth}{0pt}
\fancyfoot[C]{\small--~\thepage~--}

% \info{Label}{value} -- one line of the header block
\newcommand{\info}[2]{\textbf{#1:}~#2\par}
% \SectionLine{day and time}{room} -- one discussion section, indented
\newcommand{\SectionLine}[2]{\par\hspace*{1.5em}#1\quad[#2]}
\newcommand{\ul}[1]{\underline{#1}}

\begin{document}

%------------------------------------------------------------------------------
%  TITLE + CONTACT BLOCK
%------------------------------------------------------------------------------
\begin{center}
  {\LARGE\bfseries\color{accent}\CourseNumber}\\[0.35em]
  {\LARGE\bfseries \CourseTitle}\\[0.6em]
  {\large \Semester}\\[0.2em]
  {\large Prof.\ Alam}
\end{center}
\vspace{0.6em}
\hrule
\vspace{0.9em}

\info{Instructor}{\InstructorName}
Email: \href{mailto:\InstructorEmail}{\InstructorEmail}\par
\ul{Lectures}: \LectureRoom{} on \LectureTimes\par
Office Hours: \OfficeRoom, \OfficeHours\par

\vspace{0.6em}
\info{Teaching Assistant (TA)}{\TAName}
Email: \href{mailto:\TAEmail}{\TAEmail}\par
\ul{Lab/discussion sessions}:\TASessions\par
Office Hours: \TAOfficeRoom{} on \TAOfficeHours\par

\if\ShowExamDates1
\vspace{0.9em}
\begin{center}
\begin{tabularx}{0.85\textwidth}{@{}lX@{}}
\toprule
\textbf{Exam} & \textbf{Date} \\
\midrule
Midterm 1 & \MidtermOneDate \\
Midterm 2 & \MidtermTwoDate \\
Final exam (cumulative) & \FinalDate \\
\bottomrule
\end{tabularx}
\end{center}
\fi

%------------------------------------------------------------------------------
\section{Course description}

This course is an Introduction to Econometrics and is a required course for the
Economics major. It is also highly recommended for Econ minors and in general,
anyone who is interested in using data to understand the world around them
better. Materials of ECON 160 (\textbf{pre-requisite course}) are expected to be
well-understood by students. I strongly recommend students to refresh their
concepts of Econ 160 before we deep-dive into Econometrics.

In this course, we will learn concepts and techniques of \emph{econometric
methods} in economic analysis and implement them with real data. The empirical
models we will build, and use will be guided by intuition of economic models
building foundational bridge connecting theory and empirics. Most of the course
will focus working with \emph{cross-sectional data}, with a small part towards
the end dedicated to \emph{panel data} methods. We will not cover methods
related to time-series data in this course. This course will have a strong focus
on working with data in R. Throughout the course, in parallel, we will be
implementing the concepts learnt using R.

We will begin with a week of reviewing some concepts from ECON160, working with
R, followed by introductory concepts of econometrics which initially will build
and add mathematical formality to the concepts learned in ECON160, then dive
into new econometric methods which deal with \emph{causality}.

I want to emphasize that even though we will be using R heavily, the primary
objective in this course is to learn econometric methods, while R is simply a
tool to implement them. Hence throughout the course even though I will work with
examples with you on R, I will in most cases point you towards resources for you
to learn more by yourself since coding is just like learning any
language---you will learn the best by practicing yourself. Thus, attending
discussion sections with \TAShort---where you will be involved in coding
primarily---will be equally important as attending lectures to do well in this
course.

%------------------------------------------------------------------------------
\section{Textbook and references}

For the econometric methods only, the recommended textbook the
7\textsuperscript{th} edition of \textbf{\emph{Introductory Econometrics---A
Modern Approach}} by Jeffrey Wooldridge. Older versions of the text are
acceptable as well but please be aware that chapter numbers may not correspond
exactly between different editions. You can also use the copies of the textbook
available in the library on reserve. Before delving into the depths of the
course materials cover in this course, students must go through \textbf{Math
Refresher A} (Basic Mathematical Tools) \textbf{and B} (Fundamentals of
Probability) from Wooldridge's textbook. However, this book specializes in
teaching you econometric methods and not much into causality and econometric
techniques for causality. For the causal part of the course, as a \emph{reference
textbook}, I recommend \textbf{\emph{Mastering `Metrics: The Path from Cause to
Effect}} by Angrist and Pischke.

While working with data, we will mostly use the datasets available in the
\emph{wooldridge} package in R. These data will correspond to the textbook: the
7\textsuperscript{th} edition of \textbf{\emph{Introductory Econometrics---A
Modern Approach}} by Jeffrey Wooldridge. See
\url{https://cran.r-project.org/web/packages/wooldridge/wooldridge.pdf} for
documentation on the package and associated data within. You will need to refer
to it time and again as you will work through the problem sets and practices
that we will do in class.

Finally, even though my primary objective is to teach you econometric methods,
to start your journey on the path of becoming an expert of implementing them
with real data, there are rarely ``textbooks'' that achieves everything in the
best possible way. To that end, I will follow a combination of various
open-source resources listed below that the authors have graciously allowed the
academic community to use. Citations of resources I will be drawing from:

\begin{itemize}
  \item \emph{Statistical Inference via Data Science: A ModernDive into R and
        the Tidyverse} (2\textsuperscript{nd} Edition) by Ismai, Kim and
        Valdivia. \\ Open-source available at: \url{https://moderndive.com/v2/}
  \item \emph{R for Data Science} (2\textsuperscript{nd} edition) by Hadley
        Wickham, Mine \c{C}etinkaya-Rundel, and Garrett Grolemund. \\
        Open-source available at: \url{https://r4ds.hadley.nz/}
  \item \emph{Applied Statistics with R} by David Dalpiaz. \\
        Open-source project: \url{https://github.com/daviddalpiaz/appliedstats}
  \item \emph{Introduction to Econometrics with R} by Oswald F, Viers V,
        Villedieu P, Kennedy G (2020). \\
        Open-source project: \url{https://scpoecon.github.io/ScPoEconometrics/}
  \item \emph{Data Science for Economists} by Grant McDermott. \\
        Open-source project: \url{https://github.com/uo-ec607/lectures}
  \item \emph{Using R for Introductory Econometrics} (2\textsuperscript{nd}
        edition) by Florian Heiss. \\
        Open-source available at: \url{https://www.urfie.net/}
\end{itemize}

Although this is not necessary, an advanced reading for interested students is
\textbf{\emph{Causal Inference: The Mixtape}} by Scott Cunningham. Open-source
available at: \url{https://mixtape.scunning.com/}

%------------------------------------------------------------------------------
\section{More on coding with R}

We will use the open-source programming language R to implement the methods we
learn in this course. We will write R code in an IDE (\emph{integrated
development environment}) called RStudio. Please ensure that both \textbf{R} and
\textbf{RStudio} are installed in your machine following instructions posted in
the 1\textsuperscript{st} Canvas announcement.

If you have had no experience with any object-oriented programming language like
R before, I recommend that, in the next three weeks, you start familiarizing
yourself with the basics of R, even though I will spend the first week covering
the basics. A programming language, like any other language, takes time to get
used to. I don't want any student to fall behind just because of this---ultimately,
the goal of the course is to build your knowledge and understanding of
econometrics. There are several ways to get started:

\begin{itemize}
  \item Getting started:
        \href{https://posit.co/download/rstudio-desktop/}{Download and install R}
        and
        \href{https://posit.co/download/rstudio-desktop/}{download and install
        RStudio IDE}. It is open-source and hence free to use.
        \begin{itemize}
          \item For students unable to install RStudio, they can use
                \textbf{RStudio Cloud} as an alternative. It offers the same
                functionality as the desktop version and can be accessed
                directly at \href{https://posit.cloud/}{Posit Cloud}. However, I
                strongly recommend having it up and running on your own machine,
                since access to RStudio Cloud will always require a stable
                internet connection and is only free for 25 hours/month.
        \end{itemize}
  \item Explore resources in addition to those recommended above.
  \item Practice regularly: even just an hour a day will help you feel more
        comfortable with the basics.
  \item Join study groups: collaborating with peers can make the learning
        process easier and more enjoyable.
\end{itemize}

Working with R or any programming language will be difficult initially with a
steep learning curve, but it will pay off tremendously in the long run. This is
not only because you will learn coding but the ability to implement complicated
methods improves analytic and abstract thinking which is a very valuable in the
labor market.

\subsection{To make the most out of the course, I recommend}
\begin{itemize}
  \item Most important is to ask questions whenever you are in doubt, regardless
        of what you think about the quality of your question.
        \begin{itemize}
          \item There is no way you can learn unless you are willing to make
                mistakes.
          \item I want to encourage you to \textbf{make mistakes early in class
                by asking questions}---else the probability of mistakes
                increases in exams.
        \end{itemize}
  \item Practicing data analysis in R regularly to build your skills.
  \item Keeping up with the readings and exercises to reinforce the concepts we
        cover in class.
\end{itemize}

%------------------------------------------------------------------------------
\section{Problem sets}

Problem sets will be assigned most weeks. \TAShort{} will not accept late problem
sets (see more on this under the ``Policies'' section).

\subsection{Working format}
All assignments will include \textbf{both} theoretical calculations, coding
exercises requiring you to analyze data sets, and answer other questions based on
intuition or on the underlying math of the econometric methods that you learn. To
do all of them together in one single file, you will submit your solutions coded
and typed up in an \emph{R-notebook} file (\texttt{.Rmd}). Rmd files are designed
to allow for code, the output of code and user typed text all in one file. It
lets you develop code and record your thoughts by tightly integrating prose and
code!

\subsection{Submission on Canvas}
As you will see in the first week, R-notebooks by themselves are hard to grade
because the grader will need to download and run it on their own machine. So,
after your R-notebook with your code and typed up answers are ready, you will
``\emph{knit}'' it to a PDF which you will submit electronically on Canvas.
Canvas will \ul{not} allow you to submit any other file format. To knit to PDF
you will first need to install a \LaTeX{} distribution for your machine: MacTeX
for macOS or MiKTeX for Windows. If knitting to PDF fails, you can knit to HTML
and submit. At times installing Java may be necessary depending on your machine.
\TAShort{} will discuss R-notebooks and ``knitting'' them to PDF in the first
week in more details.

\subsection{Expectation}
Please make sure you describe \textbf{\ul{all steps}} you took to get to the
answers. You will \ul{not} receive full credit if you only give the final answer
(even if it's correct). You can discuss the problem set with your classmates, but
you need to code up your own file. If you're working with others, you must write
their names in the authors section of your R-notebook.

\subsection{Clarification/doubts on PS}
If you have any questions about the problem sets, please reach out to \TAShort{}
directly. \TAShort{} will be grading all the problem sets, and I want to ensure
consistency in the guidance provided to all students on problem sets. By
directing your problem set questions to \TAShort{} instead of me, we can avoid
any potential differences in the advice given to students. Again, since
\TAShort{} will grade, \TAthey{} will be keeping track of any late assignments.
So, if you need an extension due to unavoidable circumstances reach out to
\TAShort.

%------------------------------------------------------------------------------
\section{Grading}

\begin{itemize}
  \item \NumProblemSets{} group problem sets --- \WtProblemSets\%
        \begin{itemize}
          \item A group is a group of 2 students, no more. \emph{This is to
                foster learning from mistakes of one another and hold each other
                accountable.}
                \begin{itemize}
                  \item Inform \TAShort{} about your group by the end of the week.
                \end{itemize}
          \item One question on each exam will be randomly drawn from any of the
                PS questions.
        \end{itemize}
  \item Best of \NumQuizzesCounted{} (out of \NumQuizzes) pop-quizzes ---
        \WtQuizzes\%
        \begin{itemize}
          \item Pop-quiz will be taken in the first, or last 15 mins of any
                class. You could miss the quiz if you come late to the class.
        \end{itemize}
  \item Class attendance and participation --- \WtClassAtt\%
        \begin{itemize}
          \item Attendance sheet will only be available till the first 2 mins of
                class. Thus arriving 2 mins after the class starts will not count
                towards your attendance or class participation.
        \end{itemize}
  \item Discussion session attendance --- \WtDiscAtt\%
        \begin{itemize}
          \item Attendance sheet will only be available till the first 2 mins of
                the discussion session. Thus arriving 2 mins after the session
                starts will not count towards your attendance or class
                participation.
        \end{itemize}
  \item Participation will be both voluntary as well as randomly drawn.
  \item Midterm exam 1 (\MidtermOneDate) --- \WtMidtermOne\%
  \item Midterm exam 2 (\MidtermTwoDate) --- \WtMidtermTwo\%
  \item Final exam (\FinalDate) --- \WtFinal\% (cumulative)
  \item No exceptions and no re-take of exams allowed.
\end{itemize}

\textbf{\ul{No} collaboration will be allowed on exams.}

\subsection{Overall class grade algorithm}\hspace{0pt}\par
\vspace{0.3em}
\begin{tabular}{@{}ll@{\hspace{3em}}ll@{}}
\toprule
A  & 95--100\%   & C+ & 78--79.9\% \\
A- & 90--94.9\%  & C  & 75--77.9\% \\
B+ & 88--89.9\%  & C- & 70--74.9\% \\
B  & 85--87.9\%  & D  & 60--69.9\% \\
B- & 80--84.9\%  & F  & $<$60\% \\
\bottomrule
\end{tabular}

%------------------------------------------------------------------------------
\section{Engaged Learning Hours}

Every one unit/free credit hour course at Clark University has 180 hours of
associated engaged academic time. In this course, for the typical student, they
may be divided as follows:

\vspace{0.3em}
\begin{tabular}{@{}ll@{}}
\toprule
Class time            & 42 hours (14 weeks $\times$ 3 hours) \\
Discussion session    & 12 hours (12 classes $\times$ 1 hour) \\
Reading and class prep& 56 hours (14 weeks $\times$ 4 hours) \\
Problem sets          & 40 hours (8 $\times$ 5 hours) \\
Studying for exams    & 30 hours (3 $\times$ 10 hours) \\
\midrule
\textbf{Total}        & \textbf{180 hours} \\
\bottomrule
\end{tabular}

%------------------------------------------------------------------------------
\section{LEEP learning objectives}

As part of the Economics curriculum, this course is designed to focus on these
learning outcomes: 1) \emph{Knowledge of society}: Economics focuses on the use
of quantitative data to understand the economic aspects of society. This course
helps you think about how those data are generated and provides you with tools to
summarize and synthesize both the description of society and identify potential
associations between social phenomena (for example, between education and
income). 2) \emph{Intellectual and practical skills}: This course is all about
data: learning rigorous methods to distinguish correlation from causation and
thinking about what is ``real'' (the data generating processes) and what is due
to chance and then using statistical methods to draw conclusions from it.

%------------------------------------------------------------------------------
\section{Policies}

\subsection{No devices in class}\hspace{0pt}\par
There is enough empirical evidence that use of laptops and phones in class are
distracting to not only the student but also other students. Moreover, most of
this class will prepare you to think through econometric problems through
mathematical proofs which is best learnt when you write them yourself using pen
and paper to improve efficiency of learning. \TAShort{} will walk you through R
in the first two weeks in discussion labs where you are required to bring your
laptop. However, following that please follow \TAShort's instructions on use of
laptops in discussion classes.

This applies to AI-enabled wearables and devices as well: smart glasses (for
example, Meta Ray-Bans), AI pins and assistants, and anything else you wear or
carry that can see, hear, or answer questions for you. These must be put away and
switched off for the whole class. \textbf{During an exam, having any such device
on you counts as unauthorized assistance and will be treated as a violation of
the academic integrity policy.}

\subsection{AI usage}\hspace{0pt}\par
There is causal evidence (which we will discuss in class) that shows that
``over'' usage of AI causes students to perform worse on exams while being
stellar on their assignments. Clearly this can and will be linked to worse labor
market outcomes. So, even though we cannot monitor your AI usage, I believe that
you will minimize your AI usage for your own interest to maximize your labor
market outcomes. Given this, with advent of AI the value of critical thinking is
going to be probably at its highest. So, what I do not want is that you copy
paste your assignment into an LLM and have it produce the solution for you.
Submitting answers generated by AI without your own understanding constitutes
academic dishonesty.

\begin{enumerate}
  \item If you completely depend on AI, then you will lose the opportunity to
        build your skill and knowledge and you will prove to yourself that you
        are replaceable by AI.
  \item In the exam, you will be at a loss unless you have spent enough time
        thinking about the problems in the assignments by failing to understand
        the context of the questions.
\end{enumerate}

Students can \textbf{efficiently use AI to \ul{learn a coding language, but not
econometrics}}. The word \textbf{efficient} is important here. Efficient use of
AI is \emph{mostly} using it only in learning to code in R, by asking it to
assist you in learning the language, and not to solve your problem sets. This
requires you to be proactive and start working on your assignments as early as
possible, so that you do not have to give in to the temptation of the easy way
out of having the AI solve your assignment thereby worsening your future labor
market outcomes.

\subsection{Late work policy}\hspace{0pt}\par
Please note that all assignments should be uploaded to Canvas \textbf{before} the
due date. If you need an extension for an assignment you need to email \TAShort{}
(\href{mailto:\TAEmail}{\TAEmail}) \textbf{and ask for an extension \ul{at least
48 hours before} the assignment is due}. Once the 48-hour timeline before the
assignment is due is past, \textbf{no requests} will be accepted. This will
ensure that you can plan to work on the assignments in advance and allow
\TAShort{} to post the solutions promptly on Canvas. If you anticipate
\textbf{substantial} difficulty meeting a deadline, please reach out to me in
advance to request the problem set early, allowing you additional time to
complete it.

\subsection{Objective specific emails for efficiency / email policy}\hspace{0pt}\par
\vspace{0.4em}
\begin{center}
\begin{tabularx}{0.8\textwidth}{@{}Xl@{}}
\toprule
\textbf{Task} & \textbf{Whom to email} \\
\midrule
Extension requests on PS & \TAShort \\
Course content           & \InstructorShort{} / \TAShort{} / both \\
Question on the PS       & \TAShort \\
Coding questions         & \TAShort \\
\bottomrule
\end{tabularx}
\end{center}

\subsection{Academic integrity}\hspace{0pt}\par
All students are expected to adhere to Clark's standards of academic integrity;
this means that all work must be entirely your own and entirely unique to this
course. Plagiarism and other forms of cheating will not be tolerated or excused.
For more information, please refer to the university's policy on this issue,
available at
\url{https://catalog.clarku.edu/content.php?catoid=32&navoid=2735#academic-integrity}
or in the student handbook. If you have any questions about proper citation or
other related issues, please don't hesitate to come see me.

\subsection{Students with disabilities}\hspace{0pt}\par
Clark University is committed to providing students with documented disabilities
equal access to all university programs and facilities. Students are encouraged
to register with Student Accessibility Services (SAS) to explore and access
accommodations that may support their success in their coursework. SAS is located
on the second floor of the Shaich Family Alumni and Student Engagement Center
(ASEC). Please contact SAS at
\href{mailto:accessibilityservices@clarku.edu}{accessibilityservices@clarku.edu}
with questions or to initiate the registration process. For additional
information, please visit the SAS website at:
\url{https://www.clarku.edu/offices/student-accessibility-services/}

\subsection{Title IX}\hspace{0pt}\par
Clark University and its faculty are committed to creating a safe and open
learning environment for all students. Clark University encourages all members of
the community to seek support and report incidents of sexual harassment to the
Title IX office (\href{mailto:titleix@clarku.edu}{titleix@clarku.edu}). If you or
someone you know has experienced any sexual harassment, including sexual assault,
dating or domestic violence, or stalking, help and support is available.

Please be aware that all Clark University faculty and teaching assistants are
considered responsible employees, which means that if you tell me about a
situation involving the aforementioned offenses, I must share that information
with the Title IX Coordinator, Brittany Brickman
(\href{mailto:titleix@clarku.edu}{titleix@clarku.edu}). Although I have to make
that notification, you will, for the most part, control how your case will be
handled, including whether or not you wish to pursue a formal complaint. Our goal
is to make sure you are aware of the range of options available to you and have
access to the resources you need.

If you wish to speak to a confidential resource who does not have this reporting
responsibility, you can contact Clark's Center for Counseling and Professional
Growth (508-793-7678), Clark's Health Center (508-793-7467), or confidential
resource providers on campus: Prof.\ Stewart
(\href{mailto:als.confidential@clarku.edu}{als.confidential@clarku.edu}), Prof.\
Palm Reed
(\href{mailto:kpr.confidential@clarku.edu}{kpr.confidential@clarku.edu}), and
Prof.\ Cordova
(\href{mailto:jvc.confidential@clarku.edu}{jvc.confidential@clarku.edu}).

\subsection{FERPA}\hspace{0pt}\par
The link to Clark's policy regarding student privacy under the Family Education
Rights and Privacy Act is available here:
\url{https://www.clarku.edu/offices/security-and-identification-protection/ferpa/}

\vspace{1em}
\textbf{Disclaimer:} The instructor reserves the right to make changes to any
information contained in this syllabus at any time during the semester. Changes
will be announced, and an updated version of the syllabus will be posted on
Canvas and/or distributed to students.

\end{document}
